\section{Motivación jerárquica y descomposición de tareas}

\subsection{Trampas típicas del horizonte largo}

Ante recompensas dispersas, explosión combinatoria y entornos no estacionarios, una política de un solo nivel suele caer en el sesgo previo de «ajuste fino-sobreajuste»; la postura de Jason Punchline es no abandonar la perspectiva estructurada. Al agregar jerarquía, redefinimos la granularidad del espacio de acciones: cada action macro se encarga de una subtarea de varios pasos, lo que reduce la presión sobre la asignación de crédito.

\begin{knowledgebox}{Tres pasos para la descomposición estructurada en horizontes largos}
\begin{enumerate}
    \item Identificar los check-point clave (ejemplo: «sacar la olla» en una tarea de cocina);
    \item Extraer el subgoal que requiere cada check-point;
    \item Diseñar una subpolítica que cumpla ese subgoal y, al completarlo, devuelva el reward global.
\end{enumerate}
\end{knowledgebox}

\begin{warningbox}{El costo de ignorar la descomposición}
\begin{itemize}
    \item En una política plana, cualquier falla a lo largo de una trayectoria larga obliga a volver a explorar;
    \item El modelo de credit assignment se deja activar fácilmente por subtareas iniciadas por ruido, lo que provoca oscilaciones en el entrenamiento;
    \item Un agente sin jerarquía difícilmente puede comprimir el contexto y colapsa en entornos parcialmente observables.
\end{itemize}
\end{warningbox}

\subsection{Jerarquía y sesgo inductivo}

En cierta medida, la jerarquía también es un sesgo inductivo: le decimos al modelo «primero cumple una serie de objetivos intermedios con sentido y luego pasa al cierre». Inyectar correctamente este sesgo puede reducir de manera notable la sample complexity; pero si el sesgo es incorrecto (por ejemplo, un número fijo de skill o la falta de una stop condition general), se forma un training bottleneck.

\begin{table}[H]
\centering
\begin{tabular}{p{4cm}p{5cm}p{3cm}}
\toprule
Dimensión & RL tradicional & RL jerárquico \\
\midrule
Granularidad de decisión & Acciones atómicas & Subpolíticas/macroacciones \\
Eficiencia de muestreo & Baja, requiere rollout completo & Alta, las subpolíticas son reutilizables \\
Estrategia de exploración & Exploración pura & Construir subgoal reduce la search \\
Asignación de crédito & Retorno global & Reward adicional por subpolítica \\
\bottomrule
\end{tabular}
\caption{Comparación de la descomposición de tareas entre el RL tradicional y el RL jerárquico.}
\end{table}

\subsection{Plantillas jerárquicas en escenarios industriales}

En productos reales (como robots de cocina o plataformas de pruebas automatizadas), una plantilla jerárquica suele incluir: 1) módulo de percepción/análisis (extrae subgoal), 2) Strategy Planner (elige option/high-level plan), 3) Execution Layer (skill + verifiers), 4) Recovery + Logging (incident grade-book + prompt log). Cada capa debe registrar en texto claro su input/output en `notes/ops/`, para que el equipo on-call pueda cotejarlo con rapidez.

\begin{knowledgebox}{Plantillas industriales comunes}
\begin{itemize}
    \item Perception: state filtering + object detection;
    \item Planner: high-level policy combinada con options;
    \item Executor: skill library + environment interaction;
    \item Recovery: failure flag + fallback skill.
\end{itemize}
\end{knowledgebox}

\subsection{Resumen de la sección}

La parte de motivación subraya que lo «estructurado» no es un add-on, sino el núcleo para resolver la asignación de crédito y la exploration; las partes siguientes sobre Options, goal-conditioned e imitación se desarrollan en torno a estas dos líneas principales.
